\section{Introduction}

The weight of a coset of a linear code is the minimum Hamming weight of a
vector in that coset.  We determine the cosets of weight at least \(r-1\) for
a family of long generalized and extended Reed--Solomon codes of redundancy
\(r\).  The same classification determines their deep holes and the
one-column extensions whose Singleton defect is at most one.

Let \(A\subseteq\PP^1(\F_q)\), put \(s=|A|\) and
\(S=\PP^1(\F_q)\setminus A\), and let \(C_S(v)\) be a redundancy-\(r\)
generalized Reed--Solomon code on \(S\), with arbitrary nonzero coordinate
multipliers \(v\).  A parity-check matrix \(H_S\) has columns proportional to
the points of the degree-\((r-1)\) normal rational curve
\[
  \NRC_{r-1}=\{(1,t,\ldots,t^{r-1}):t\in\F_q\}
              \cup\{(0,\ldots,0,1)\}\subset\PP^{r-1}.
\]
For a syndrome \(f\), let \(d_S(f)\) be the least number of columns of \(H_S\)
whose span contains \(f\); this is the weight of the corresponding coset.
Column multipliers do not change these projective spans.

In characteristic two, at redundancies six and seven only, a second linear
locus appears: adjacent zeros in Pascal row \(r-2\) determine what we call
the maximal Lucas carrier.

\begin{theorem}[High-weight cosets]
\label{thm:main}%
\coverage{absent}%
\uses{thm:simultaneous-marker-escape,prop:trapped-row-space}%
\imports{duer,seroussi-roth,zwk-full-support-rank-two}%
\proves{main-fixed-level-classification-package}%
Let \(r\geq6\), let \(A\subseteq\PP^1(\F_q)\) have size \(s\), let
\(p=\operatorname{char}\F_q\), and assume \(r\leq q+1-s\).  Put
\[
 Q^*_{r,s}=6(r+s)-16+
 \left\lfloor2\sqrt{6(r+s)-18}\right\rfloor .
\]
Let \(\cP_r\) be the rank-at-most-two locus of the consecutive
three-row Hankel (second catalecticant) matrix of \(f\), and let
\(\cE_{r,p}=\cP_r\), except that
\(\cE_{r,2}=\cP_r\cup\cM^{\max}_{r,2}\) for \(r\in\{6,7\}\), where
\[
 \cM^{\max}_{r,p}=
 \PP\!\left\langle e_j:
 0\leq j\leq r-1,\quad
 \binom{r-2}{j}=\binom{r-2}{j-1}=0\pmod p\right\rangle
\]
in the standard divided-power coordinate basis \(e_0,\ldots,e_{r-1}\), with
\(\binom nk=0\) for \(k\notin\{0,\ldots,n\}\); explicitly
\(\cM^{\max}_{6,2}=\PP\langle e_2,e_3\rangle\) and
\(\cM^{\max}_{7,2}=\{[e_3]\}\).
If \(q\geq Q^*_{r,s}\), then
\[
 f\notin \cE_{r,p}
 \quad\Longrightarrow\quad d_S(f)\leq r-2.
 \autoeqtag{eq:main-carrier-containment}
\]
In characteristic two the same conclusion holds with
\(6(r+s)-22+\lfloor2\sqrt{6(r+s)-24}\rfloor\) in place of
\(Q^*_{r,s}\).

Assume in addition that \(p\) is odd, or that \(p=2\) and \(r\geq8\), so
that \(\cE_{r,p}=\cP_r\).  If \(s=0\), the covering radius is
\(r-1\), and the cosets of weight \(r-1\) are exactly the tangent and
conjugate-secant directions.  If \(s>0\), the covering radius is \(r\): the
weight-\(r\) directions are the omitted curve points, and the weight-\(r-1\)
directions are exactly the tangent points, conjugate-secant points, and
interior points of rational secants incident with \(A\).  Their projective
counts are
\[
 N_r=s,\qquad
 N_{r-1}=\frac{q(q+1)^2}{2}
       +\frac{(q-1)s(2q+1-s)}2.
\]
The corresponding numbers of nonzero cosets are \((q-1)N_r\) and
\((q-1)N_{r-1}\).
\end{theorem}

Geometrically, \(\cP_r\) is the second secant variety of
\(\NRC_{r-1}\).  Its \(\F_q\)-points comprise the curve itself and the
rational-secant, tangent, and conjugate-secant strata.

Here \(Q^*_{r,s}=6(r+s)+O(\sqrt{r+s})\).  The classification has no
characteristic hypothesis beyond the two binary cases \(r\in\{6,7\}\).
For full support, Theorems~\ref{thm:r6} and~\ref{thm:r7} settle those over
every prime power \(q\geq7\): the carrier contributes one nucleus orbit,
respectively one central nucleus point, exactly when \(q=2^m\) with \(m\)
odd; the exclusion is therefore a phenomenon, not a limitation of the
method.  For \(s>0\) in those two cases, \eqref{eq:main-carrier-containment}
still holds and only the arithmetic of the carrier is left open.  The Lucas carrier \(\cM^{\max}_{r,p}\) is nonempty for many other
pairs \((r,p)\) and is a genuine component of the recursively contained
carrier of Theorem~\ref{prop:recursive-component-selection}, but a
branchwise reading of that theorem shows that it traps no syndrome above
the threshold (Proposition~\ref{prop:trapped-row-space}).  Deleting
even one point is substantive: the radius jumps from \(r-1\) to \(r\), and the
omitted points become exactly the new weight-\(r\) shell.

The only computer-algebra input to the uniform theorem is the exact reduced
terminal-carrier decomposition in Proposition~\ref{prop:reduced-terminal-carrier},
verified by the registered Gr\"obner-elimination certificate.  Recursive
component selection, simultaneous finite-field selection, shell transfer, and
the enumerative double counts are mathematical arguments.  The covering-radius
and MDS-extension gates are imported from Seroussi--Roth and D\"ur.

\noindent\textbf{Coding interpretation.}
Theorem~\ref{thm:main} classifies every coset in the two highest possible
weight shells of these GRS codes and, equivalently, every one-column extension
having Singleton defect at most one.
Indeed, appending \(f\) as a parity-check column gives a code \(\widehat C_f\)
with \(d(\widehat C_f)=d_S(f)+1\).  A family-wise double count
determines the aggregate number of minimum words in the tangent,
conjugate-secant, and incident-split-secant families; the standard NMDS
recurrence then gives their complete family-aggregate weight enumerators.

For \(s=0\), \(C_S(v)\) is monomially equivalent to the full extended or
projective Reed--Solomon code.  For \(s>0\), a projective change of coordinate moves an
omitted point to infinity, so the code is monomially equivalent to an
ordinary GRS code.  To our knowledge, no earlier result classifies, for
arbitrary redundancy, every projective syndrome direction of such a code
with coset weight at least \(r-1\) after finitely many prescribed points are
deleted.  In the range of Theorem~\ref{thm:main}, these are exactly the
omitted curve points in weight \(r\), and the tangent, conjugate-secant, and
deleted-point-incident split-secant directions in weight \(r-1\).

The qualification is essential.  The covering-radius, syndrome, and
MDS-extension dictionaries are classical
\cite{SeroussiRoth1986,Dur1994,Kaipa2017,WuDingChen2023}.  For full projective
support at redundancy four, the tangent and conjugate-secant shell and its
count are already known
\cite{ZWK2020,BPS2022,DMPCosets2021}.  Point-deleted projective evaluation
sets were studied in \cite{XuHongXu2017}, and recent work gives general
MDS/NMDS constructions from deep holes \cite{LiLuLingLam2026}.  None of these
results gives the complete point-deleted top two shells at arbitrary
redundancy.  The new contribution is the
syndrome-to-Hankel carrier classification together with simultaneous
exclusion off that carrier; these yield the arbitrary-redundancy
prescribed-deletion shell theorem.  The extension and enumerator statements
are consequences.

The witness sought is a degree-\((r-2)\) binary form in \(W_f\) whose distinct
rational roots all lie in \(S\); by the Hankel criterion, its roots are the
\(r-2\) parity-check columns spanning \(f\).  Instead of classifying every
intermediate redundancy, the proof selects \(r-5\) of these linear factors
simultaneously.  Contraction leaves a redundancy-five (R5) pencil of binary
cubics, and any split squarefree member avoiding the selected roots and \(A\)
completes the desired form.  The carrier theorem confines the cases in which
every such contraction is trapped to the persistent rank-two locus \(\cP_r\),
together with the line \(\PP\langle e_2,e_3\rangle\) at \((r,p)=(6,2)\)
and the point \([e_3]\) at \((7,2)\); outside that locus, a
degree-six selector, a Vandermonde choice, and the exact R5 count produce the
cubic.  The witnesses are abundant: there are at least
\(q^{r-4}/(r-2)!-O_{r,s}(q^{r-9/2})\) split locators supported on \(S\).

The full-support calculations remain useful where they are sharper than the
uniform theorem.  Redundancy five supplies the terminal cubic engine and a
complete classification over every prime power \(q\geq7\); redundancies six
and seven supply exact small-field and modular refinements.  They are stated
and proved after the arbitrary-redundancy theorem, with their separate radius
and certificate boundaries kept explicit.  The R8--R10 calculations are not
used as induction hypotheses and are not included here.
